\documentclass[10pt,a4paper]{article}

%% pacotes essenciais e atualizados
\usepackage[T1]{fontenc}      % Essencial para hifenização de palavras acentuadas
\usepackage[utf8]{inputenc}
\usepackage[brazilian]{babel}    % Tradução de termos automáticos e regras pt-BR
\usepackage[top=2.5cm, bottom=2.5cm, left=2cm, right=2cm]{geometry} % Controle de margens
% Adicionado o pacote amssymb para usar o comando \square
\usepackage{graphicx,color,verbatim,Sweave,url,xargs,amsmath,amssymb,booktabs,longtable,eurosym}

%% novos ambientes requeridos pelo pacote exams
\newenvironment{question}{\item}{}
\newenvironment{solution}{\comment}{\endcomment}
\newenvironment{answerlist}{\renewcommand{\labelenumii}{(\alph{enumii})}\begin{enumerate}}{\end{enumerate}}

%% compatibilidade com pandoc e knitr/rmarkdown
\providecommand{\tightlist}{\setlength{\itemsep}{0pt}\setlength{\parskip}{0pt}}
\providecommand{\pandocbounded}[1]{#1}
\setkeys{Gin}{keepaspectratio}

%% fontes: Helvetica (sem serifa, excelente para provas)
\usepackage{helvet}
\IfFileExists{sfmath.sty}{
  \RequirePackage[helvet]{sfmath}
}{}
\renewcommand{\sfdefault}{phv}
\renewcommand{\rmdefault}{phv}

%% novos comandos para receber variáveis do R
\makeatletter
\newcommand{\Course}[1]{\def\@Course{#1}}
\newcommand{\Semester}[1]{\def\@Semester{#1}}
\newcommand{\ID}[1]{\def\@ID{#1}}

\Course{Disciplina}
\Semester{Ano/Semestre}
\ID{00000}

%% \exinput{header}

\newcommand{\myCourse}{\@Course}
\newcommand{\mySemester}{\@Semester}
\newcommand{\myID}{\@ID}
\makeatother

%% cabeçalhos das páginas secundárias
\markboth{\textnormal{\bf \large \myCourse \: \mySemester \hfill \myID}}%
         {\textnormal{\bf \large \myCourse \: \mySemester \: \myID}}
\pagestyle{myheadings}

\begin{document}

%% página de rosto / cabeçalho principal
\thispagestyle{empty}
{\sf
\begin{center}
    \textbf{\LARGE Universidade Federal de Santa Catarina}\\[0.15cm]
    \textbf{\Large Campus de Curitibanos}\\[0.5cm]
    \textbf{Exercício de \large \myCourse \ -- Semestre \mySemester} \hfill \textbf{ID: \myID}
\end{center}

\vspace{0.3cm}

% Quadro de identificação do aluno e espaço para correção
\noindent\fbox{
  \begin{minipage}{\textwidth}
    \vspace{0.2cm}
    \textbf{Nome:} \hrulefill \\[0.4cm]
    \textbf{Matrícula:} \rule{3.5cm}{0.4pt} \hfill
    \textbf{Data:} \rule{2.5cm}{0.4pt} \hfill
    \textbf{Avaliação:} $\square$ Correto \hspace{0.2cm} $\square$ Incorreto
    \vspace{0.2cm}
  \end{minipage}
}
}

\vspace{0.2cm}

% Aviso importante sobre a resolução
\vspace{0.5cm}
\begin{center}
    \fbox{
    \begin{minipage}{0.9\textwidth}
        \centering
        \textbf{\textit{Atenção: Apresente uma resolução completa e didática, explicitando todos os passos e o raciocínio utilizado. Estruture sua resposta para que sirva como material de estudo futuro. Respostas diretas sem o devido desenvolvimento não serão aceitas.}}
    \end{minipage}
    }
\end{center}
\vspace*{0.5cm}

% Início dos exercícios
\begin{enumerate}
\setlength{\itemsep}{1.5em} % Aumenta levemente o espaço entre uma questão e outra
%% \exinput{exercises}
\end{enumerate}

\end{document}
